% Ejemplo ilustrativo (cifras de resultados hipotéticas). Sustitúyelo por el texto de tu investigación.

\guia{Las conclusiones representan entre 5 y 8\% del documento: no son un resumen, son el cierre argumentativo. Retoman el problema, confirman el objetivo y la hipótesis y cierran con la aportación.}

\etiqueta{1. Retoma del problema}
Esta investigación partió de una paradoja: aunque el peso mexicano es una de las monedas emergentes más negociadas del mundo y enfrentó episodios de alta volatilidad, el volumen de futuros negociado en el MexDer se redujo alrededor de 86\% entre 2009 y 2022. Este hecho cuestiona la idea de que la demanda de cobertura cambiaria se canaliza de forma natural hacia el mercado local de derivados.

\etiqueta{2. Confirmación del objetivo}
El objetivo de analizar la relación entre la volatilidad del tipo de cambio peso-dólar y el volumen de negociación de futuros del dólar en el MexDer se alcanzó. La combinación de un modelo GARCH(1,1) con regresiones por subperíodos permitió distinguir cómo cambió esa relación antes y después de 2016.

\etiqueta{3. Comprobación de la hipótesis}
La hipótesis se confirmó: entre 2009 y 2015, un aumento de un punto porcentual en la volatilidad condicional se asoció con un incremento de 4.2\% en el volumen negociado ($\beta = 0.042$, $p < 0.01$), mientras que entre 2016 y 2022 el efecto dejó de ser significativo ($\beta = 0.011$, $p = 0.21$).

\etiqueta{4. Hallazgos principales}
Primero, la sensibilidad del volumen del MexDer a la volatilidad cambiaria se redujo a una cuarta parte en el segundo subperíodo. Segundo, el volumen de futuros del peso en el CME se relaciona negativamente con la actividad del MexDer ($\beta = -0.351$, $p < 0.01$), lo que es consistente con una sustitución entre mercados. Tercero, el diferencial de tasas de interés perdió significancia después de 2016, lo que sugiere que las estrategias de carry trade dejaron de canalizarse por el mercado local.

\etiqueta{5. Limitaciones}
El estudio presenta dos limitaciones principales: la falta de información mensual sobre la cobertura cambiaria pactada en el mercado extrabursátil local y el uso de datos agregados, que no permiten identificar el tipo de participante (cobertura o especulación) detrás de cada contrato.

\etiqueta{6. Líneas futuras de investigación}
Se recomienda analizar, con datos por tipo de participante, si la migración hacia el CME responde a costos de transacción, a requisitos de margen o al acceso de inversionistas extranjeros, así como extender el análisis a los contratos de futuros sobre la TIIE y sobre el Índice de Precios y Cotizaciones.

\etiqueta{7. Aportación final — el broche de oro}
La aportación central de esta investigación consiste en mostrar que el problema del MexDer no es la falta de demanda de cobertura, sino su incapacidad para retener la liquidez frente a mercados externos y extrabursátiles. Esta evidencia reorienta el debate de política financiera: más que promover el uso de derivados, la prioridad es reducir los costos y las barreras que hacen menos competitivo al mercado local.
